\section{Cohort Audit, Outcome Reconstruction, and Leakage Removal}

\runin{Provenance and Skip-Logic Reconstruction}
All 368,132 isolated-CABG procedures recorded in STS ACSD v2.81 (July 1, 2014--June 30, 2017; 1,091 U.S.\ centers across all nine U.S.\ Census divisions) were pinned cryptographically prior to modeling. Under STS ACSD v2.81 Sequence 4930 (\cli{COtAFib}), the atrial fibrillation complication indicator is populated only when the parent complication gate \cli{Complics} equals 1 (any postoperative complication).\cite{sts2014v281} Among 232,783 uncomplicated recoveries (\cli{Complics == 2}, 63.31\% of records), \cli{COtAFib} is structurally blank by skip logic rather than missing at random. Treating structural blanks as missing introduces severe collider bias by conditioning on postoperative morbidity.\cite{hernan2020causal} Resolving \cli{Complics == 2} and (\cli{Complics == 1}, \cli{COtAFib == 2}) as POAF-negative and (\cli{Complics == 1}, \cli{COtAFib == 1}) as POAF-positive recovered \textbf{367,695 labeled cases} (437 genuinely unknown excluded; \textbf{89,337 POAF events, 24.30\%}). Clinical validity was confirmed by expected risk separation between positive and negative cases (mean age 68.5 vs.\ 63.2\,yr; discharge amiodarone 74.3\% vs.\ 14.7\%).

\runin{Strict 638 $\to$ 207 Pre-Incision Variable Contract}
Exhaustive column-by-column auditing against STS ACSD v2.81 specifications eliminated 431 ineligible fields, including seven early-column postoperative duration leaks (\cli{LOS}, \cli{PLOS}, \cli{ORDuration}, \cli{SIDuration}, \cli{SurgtoMort}, \cli{SurgtoRead}, \cli{SurgtoSternInf}), all intraoperative bypass/cross-clamp durations, blood products, discharge medications, STS risk calculators, auxiliary postoperative outcomes, and hospital/surgeon identifiers. Exactly \textbf{207 raw pre-incision predictors} remained (including pre-incision operating-room \cli{PASYS} where measured prior to skin incision, plus 9-level \cli{Region}). Sentinel non-numeric codings (e.g., 999) were mapped cleanly to missing values prior to model ingestion.

\runin{Dual Controlled Validation Splits}
To evaluate both cold-start external transport and in-network hospital deployment cleanly, we constructed two matched evaluation cohorts with identical holdout size ($N{=}21{,}241$, natural prevalence $\approx 24.3\%$) and identical 50,000-patient balanced five-fold hospital-grouped training budgets:
\begin{enumerate}
  \item \textbf{Regime 1 — Unseen Hospitals (0.0\% Overlap):} Multivariate anticlustering partitioned the 1,091 centers into 1,017 development hospitals (346,454 procedures; nested 50,000-patient balanced benchmark across five exact hospital-disjoint CV folds) and 74 strictly unseen holdout hospitals ($N{=}21{,}241$; 5,160 events; $\mathrm{Max}\,|\mathrm{SMD}|=0.103$ across 206 numeric variables).\cite{papenberg2021anticlustering,joseph2022split}
  \item \textbf{Regime 2 — Seen Hospitals (100.0\% Overlap):} Hospital-stratified partitioning guaranteed that all 1,091 surgical centers appeared in both the 346,454-case development pool and the nested 50,000-patient benchmark, while reserving a matched natural-prior holdout ($N{=}21{,}241$; 1,086 seen centers; zero unseen centers). For site-calibrated in-network analyses, hospital POAF prevalence priors were estimated strictly on training records via Empirical-Bayes shrinkage toward the national prior ($\mathrm{EB}_{50\mathrm{k}}$ with $m_0{=}15$; $\mathrm{EB}_{346\mathrm{k}}$ with $m_0{=}25$).
\end{enumerate}

\section{Extended Clinical Operating Points, Subgroup Equity, and TRIPOD+AI Checklist}

\begin{table}[H]
\centering
\caption{\ct{Decision Curve Net Benefit and Clinical Operating Characteristics for NEVERHUMAN Across 74 Unseen U.S.\ Hospitals ($N{=}21{,}241$).}}
\label{tab:supp_dca}
\tablefont
\setlength{\tabcolsep}{3.5pt}
\begin{tabularx}{\textwidth}{@{}R{1.1cm} R{1.3cm} R{1.2cm} R{1.2cm} R{1.2cm} R{1.2cm} R{1.5cm} R{1.5cm} Y@{}}
\toprule
\thd{Thresh} & \thd{Flagged} & \thd{Sens} & \thd{Spec} & \thd{PPV} & \thd{NPV} & \thd{NB (NH)} & \thd{NB (POAF)} & \thd{Avoided / 100 vs.\ Treat-All} \\
\midrule
15\% & 77.8\% & 91.5\% & 26.6\% & 28.5\% & 90.6\% & 0.1240 & 0.1161 & \textbf{+8.3} \\
20\% & 63.1\% & 81.3\% & 42.7\% & 31.3\% & 87.7\% & 0.0890 & 0.0737 & \textbf{+14.2} \\
25\% & 45.1\% & 64.5\% & 61.1\% & 34.8\% & 84.3\% & 0.0587 & 0.0407 & \textbf{+20.4} \\
30\% & 29.3\% & 45.6\% & 75.9\% & 37.8\% & 81.3\% & 0.0325 & 0.0140 & \textbf{+26.6} \\
35\% & 17.5\% & 29.8\% & 86.5\% & 41.5\% & 79.3\% & 0.0174 & 0.0021 & \textbf{+33.8} \\
40\% & 9.8\% & 18.2\% & 92.8\% & 44.8\% & 78.0\% & 0.0078 & 0.0004 & \textbf{+40.4} \\
\bottomrule
\end{tabularx}
\end{table}

\begin{table}[H]
\centering
\caption{\ct{Prespecified 25-Subgroup Equity and Clinical Transport Audit on 74 Unseen U.S.\ Hospitals ($N{=}21{,}241$; NEVERHUMAN Wins 24/25 Subgroups vs.\ CatBoost 50k).}}
\label{tab:supp_equity}
\tablefont
\setlength{\tabcolsep}{3.4pt}
\begin{tabularx}{\textwidth}{@{}L{2.3cm} Y R{1.15cm} R{1.15cm} R{1.45cm} R{1.45cm} R{1.25cm} R{1.15cm}@{}}
\toprule
\thd{Axis} & \thd{Subgroup} & \thd{$n$} & \thd{POAF \%} & \thd{NEVERHUMAN} & \thd{CatBoost 50k} & \thd{$\Delta$ (bps)} & \thd{NH ECE} \\
\midrule
Overall & All Unseen-Hospital Patients & 21,241 & 24.3\% & \textbf{0.6767} & 0.6742 & +24.5 & 0.0072 \\
Sex & Male & 16,062 & 25.1\% & \textbf{0.6719} & 0.6698 & +21.1 & 0.0085 \\
Sex & Female & 5,172 & 22.0\% & \textbf{0.6891} & 0.6854 & +37.2 & 0.0109 \\
Age & Age $<$ 60 yr & 6,283 & 13.2\% & \textbf{0.6667} & 0.6668 & $-$1.2 & 0.0074 \\
Age & Age 60--69 yr & 7,881 & 24.4\% & \textbf{0.6149} & 0.6104 & +45.4 & 0.0113 \\
Age & Age 70--79 yr & 5,627 & 33.0\% & \textbf{0.6173} & 0.6130 & +43.2 & 0.0134 \\
Age & Age $\ge$ 80 yr & 1,450 & 38.3\% & \textbf{0.5843} & 0.5758 & +84.7 & 0.0321 \\
Race / Ethnicity & White & 17,206 & 25.7\% & \textbf{0.6688} & 0.6664 & +24.2 & 0.0083 \\
Race / Ethnicity & Black & 1,822 & 17.1\% & \textbf{0.6849} & 0.6835 & +13.8 & 0.0136 \\
Race / Ethnicity & Asian & 842 & 18.8\% & \textbf{0.6534} & 0.6472 & +62.4 & 0.0317 \\
Race / Ethnicity & Hispanic / Latino & 1,433 & 16.5\% & \textbf{0.7146} & 0.7109 & +37.3 & 0.0244 \\
Comorbidity & Diabetes Mellitus & 10,456 & 24.1\% & \textbf{0.6812} & 0.6787 & +24.7 & 0.0106 \\
Comorbidity & No Diabetes & 10,738 & 24.5\% & \textbf{0.6736} & 0.6708 & +27.4 & 0.0116 \\
Comorbidity & LVEF $<$ 40\% & 2,893 & 25.8\% & \textbf{0.6787} & 0.6761 & +26.4 & 0.0231 \\
Comorbidity & LVEF $\ge$ 40\% & 17,764 & 24.1\% & \textbf{0.6764} & 0.6741 & +22.9 & 0.0073 \\
Comorbidity & Chronic Lung Disease & 4,643 & 27.0\% & \textbf{0.6710} & 0.6691 & +19.3 & 0.0175 \\
Comorbidity & CKD (eGFR $<$ 60) & 4,561 & 28.9\% & \textbf{0.6439} & 0.6417 & +22.5 & 0.0118 \\
Acuity & Elective CABG & 7,887 & 25.0\% & \textbf{0.6652} & 0.6649 & +3.3 & 0.0118 \\
Acuity & Urgent / Emergent CABG & 13,349 & 23.9\% & \textbf{0.6830} & 0.6794 & +35.8 & 0.0081 \\
Stress Slice & No Recorded Preop AF & 21,103 & 24.3\% & \textbf{0.6771} & 0.6746 & +24.6 & 0.0068 \\
Stress Slice & Rare Clinical Tails & 1,917 & 26.2\% & \textbf{0.6935} & 0.6890 & +44.7 & 0.0273 \\
Stress Slice & 2017 Unseen Hospitals & 3,555 & 22.5\% & \textbf{0.6579} & 0.6565 & +14.4 & 0.0264 \\
Hospital Volume & Low-Vol Centers (n$\le$141) & 1,864 & 26.0\% & \textbf{0.6818} & 0.6784 & +34.3 & 0.0310 \\
Hospital Volume & Mid-Vol Centers & 5,282 & 23.0\% & \textbf{0.6851} & 0.6830 & +21.4 & 0.0122 \\
Hospital Volume & High-Vol Centers (n$>$341) & 14,095 & 24.6\% & \textbf{0.6727} & 0.6702 & +24.7 & 0.0080 \\
\bottomrule
\end{tabularx}
\end{table}

\begin{table}[H]
\centering
\caption{\ct{TRIPOD+AI Adherence Summary for the National STS ACSD v2.81 Pre-Incision POAF Benchmark.}}
\label{tab:supp_tripod}
\tablefont
\setlength{\tabcolsep}{4pt}
\begin{tabularx}{\textwidth}{@{}L{3.2cm} L{2.0cm} Y@{}}
\toprule
\thd{TRIPOD+AI Domain} & \thd{Location} & \thd{Implementation Summary} \\
\midrule
Data Source \& Provenance & Methods, Table 1 & Full national STS ACSD v2.81 isolated-CABG export (2014--2017; 368,132 procedures; 1,091 U.S.\ centers; SHA-256 verified). \\
Outcome Definition & Methods, Supp.\ S1 & Skip-logic reconstruction (\cli{Complics} $\to$ \cli{COtAFib}) recovering 367,695 labeled cases (89,337 POAF events, 24.30\%). \\
Predictor Audit & Methods, Supp.\ S1 & Strict 638 $\to$ 207 pre-incision variable audit purging 7 early-column postoperative leaks and all intraoperative/postoperative fields. \\
Validation Regimes & Methods, Tables 1--2 & Dual controlled splits: Regime 1 Unseen Hospitals (74 unseen centers, 0\% overlap) and Regime 2 Seen Hospitals (all 1,091 centers in train, 100\% overlap + train-only EB hospital calibration). \\
Comparators \& Parity & Methods, Tables 2--3 & All five public comparators (TabFM v1.0.0, KumoTabular-Large v2, TabPFN-3.5, CatBoost, AutoGluon \cli{best\_quality}) run strictly on binary $\{0,1\}$ targets across identical 50,000-patient 5-fold hospital-grouped cohorts vs.\ proprietary NEVERHUMAN. \\
Calibration \& Clinical Utility & Results, Fig.\ 5, S1 & Uniform fold-OOF Platt calibration + dev-only Bayes shift; 15-bin ECE, slope, O:E, 7.4$\times$ decile spread, and Decision Curve Analysis. \\
Uncertainty \& Equity & Results, Table 2, S2 & Exact 5,000-draw whole-hospital cluster bootstrap 95\% CIs (seed 42); prespecified 25-subgroup equity audit (24/25 wins). \\
Transparency \& Governance & Discussion / Statements & Explicit disclosure of NEVERHUMAN proprietary status, retrospective selection bounds, and required frozen prospective STS v4.20+ confirmation. \\
\bottomrule
\end{tabularx}
\end{table}

